\begin{abstract}
Interventions on connectome-constrained models are often reported through the
graph statistics they preserve. Such reports can miss the task information
that reaches a readout. We study this failure for a fixed directed operator:
two rounds of incoming-neighbor averaging followed by a readout at two
designated sinks. Propagating the two source basis columns summarizes every
intervention by a $2\times2$ source-to-readout transfer $T$. For the noiseless
subspace XOR used in prior experiments, both columns being nonzero are
necessary and sufficient for perfect Bayes accuracy, while zeroing a common
relay is neither. On reconstructed MaleCNS graphs, two edits that preserve
degree and module counts and eliminate all common relays differ in exactly
this way; their archived accuracies are 0.9745 and 0.5055. Across two $k=8$
extractions, all nonzero doses keep ideal label information while finite
accuracy rises from about 67\% to 94\%, and several null graphs regain the
task after a zero-dose graph loses it. An executed MLP reproduces the
high-versus-chance split on the matched matrices, and a visual parity XOR
separates an invertible transfer (0.746), a parallel mixture (0.663), and a
deleted source (0.506). With four CIFAR-10 images mixed together, the class
at one source changes from 0.164 without a private transfer direction to
0.566 with one, and falls to 0.096 when that source is deleted. The same
graph still reads a preserved source at 0.440. The result is
a task-specific reporting standard: state the executed operator and check the
information reaching the readout before interpreting a structural control.
\end{abstract}
